

%%%%%%%%%HO SPOSTATO NEL TESTO PRINCIPALE QUESTA DIMOSTRAZIONE CHE ERA L'UNICA PRESENTE NELL'APPENDICE ALLA SEZIONE 2%%%%%%%%%%%%% 
\iffalse
\section{Appendix to Section \ref{sec:pca}}\label{app:sec:pca}

\begin{proof}[Proof of Theorem~\ref{thm:freeenriched}]%\marginpar{To be moved in the appendix}\marginpar{Questo risultato di sicuro esiste in letteratura}
The result follows from more general results about enriching on arbitrary algebraic theories: see for instance \cite[Prop. 6.4.7]{borceux2} or \cite[Cor. 1]{villoria2024enriching}. It is anyway convenient to show the involved structures.

We first illustrate the definition of $F^+\colon \Cat{C}^+ \to \Cat{D}^+$ for a functor $F\colon \Cat{C} \to \Cat{D}$. For all objects $X$, $F^+$ is defined as $F^+(X)\defeq F(X)$; For all $d\in \Cat{C}^+[X,Y]$, $F^+(d)\in \Cat{D}^+[FX, FY]$ is defined for all $h\colon FX \to FY$ as $F^+(d)(g)\defeq \sum_{\{f\in \Cat{C}[X,Y]| F(f)=g\}}d(f)$. 
One can easily check that the functor is PCA-enriched.

For all categories $\Cat{C}$, the unit $\eta_{\Cat{C}}\colon \Cat{C} \to \Cat{C}^+$ is the functor acting as identity on objects and mapping any arrow $f\in \Cat{C}[X,Y]$ into $\delta_f \in \Cat{C^+}[X,Y]$.

Now, given a PCA-enriched category $\Cat{D}$ and a functor $F\colon \Cat{C} \to U(\Cat{D})$, one can define a functor $F^\sharp \colon \Cat{C}^+ \to \Cat{D}$ as follows: for all objects $X$, $F^\sharp(X)\defeq F(X)$, for all arrow $d\in \Cat{C}^+[X,Y]$, $F^\sharp(d)\defeq \sum_{f\in \Cat{C}[X,Y]}d(f)\cdot \delta_{F(f)}$.   
\end{proof}
\fi


\section{Appendix to Section \ref{sec:pca}}

\begin{proof}[Proof of Lemma~\ref{lemma:initialproperties2}]
    \begin{enumerate}
        \item  Follows easily by the enrichment. For instance, the first equality is proved as follows.
\begin{align*} 
(p\cdot f) ; g & = (f+_p\star_{X,Y});g \tag{def of $p\cdot(-)$} \\
&= (f;g) +_p(\star_{X,Y};g) \tag{\ref{eq:enr}}\\
&= (f;g) +_p(\star_{X,Z}) \tag{\ref{eq:enr}}\\
&= p\cdot (f ; g) \tag{def.\ of $p\cdot(-)$} \\
\end{align*}
\item Follows easily by the laws of PCAs.
\begin{align*}
p\cdot (q \cdot f) &= (f+_q \star )+_p \star \tag{def.\ of $p\cdot(-)$}\\
&= f+_{pq} (\star +_{\frac{p\cdot(1-q)}{1-pq}} \star) \tag{\ref{eq:pca}} \\
&= f+_{pq} \star \tag{\ref{eq:pca}} \\
&= pq \cdot f \tag{def.\ of $p\cdot(-)$}
\end{align*}
%\item Follows from Lemma~\ref{lemma:initialproperties}.\ref{lemma:p;} and induction. The case $n=0$ follows by enrichment. For $n+1$ consider
%\begin{align}
%   ( \sum_{i=1}^{n+1}p_i\cdot f_i);h&= ( f_{n+1}+_{p_{n+1}}\sum_{i=1}^{n}p_i\cdot f_i);h \tag{Def. of $\sum_{i=1}^{n+1}{p_i}\cdot f_i$ in (\ref{eq:def somma n pca})}\\
%   &= f_{n+1};h +_{p_{n+1}} (\sum_{i=1}^{n}\frac{p_i}{1-p_{n+1}}\cdot f_i);h \tag{PCA-enrichment}\\
%   &= f_{n+1};h +_{p_{n+1}} \sum_{i=1}^{n}\frac{p_i}{1-p_{n+1}}\cdot (f_i;h) \tag{Lemma~\ref{lemma:initialproperties}.\ref{lemma:p;} + Ind. Hp.}\\
%   &= \sum_{i=1}^{n+1}p_i\cdot (f_i;h). \tag{Def. of $\sum_{i=1}^{n+1}p_i\cdot f_i$ in (\ref{eq:def somma n pca})}
%\end{align}
%The other case is symmetric.
\item Follows by induction. The case $n=0$ follows from the laws of PCA: $q\cdot \star= \star +_q \star= \star$. For $n+1$ consider 
\begin{align}
    q\cdot\sum_{i=1}^{n+1}p_i\cdot f_i &= q\cdot ( f_{1} +_{p_{1}}\sum_{j=1}^{n}q_j\cdot f_j) \tag{Def. of $\sum_{i=1}^{n+1}{p_i}\cdot f_i$ in (\ref{eq:def somma n pca})}\\
    &= q\cdot f_{1} +_{p_{1}} \sum_{j=1}^{n}qq_j\cdot f_j \tag{PCA laws + Ind. Hp.}\\
    &= \sum_{i=1}^{n+1}qp_i\cdot f_i \tag{Def. of $\sum_{i=1}^{n+1}{p_i}\cdot f_i$ in (\ref{eq:def somma n pca})}
\end{align}
%\item Follows by induction. The case $n=0$ follows by enrichment. For $n+1$ consider 
%\begin{align}
%    \sum_{i=1}^{n+1}p_i\cdot (f_i+_p g_i)&= (f_{n+1}+_p g_{n+1}) +_{p_{n+1}}\sum_{i=1}^{n}p_i\cdot (f_i+_p g_i) \tag{Def. of $\sum_{i=1}^{n+1}{p_i}\cdot f_i$ in (\ref{eq:def somma n pca})}\\
 %   &=(f_{n+1}+_p g_{n+1}) +_{p_{n+1}} ((\sum_{i=1}^{n}p_i\cdot f_i)+_p ( \sum_{i=1}^{n}p_i\cdot g_i)) \tag{Ind. Hp.}\\
 %   &= (f_{n+1}+_{p_{n+1}} (\sum_{i=1}^{n}p_i\cdot f_i)) +_{p} (g_{n+1}+_{p_{n+1}} ( \sum_{i=1}^{n}p_i\cdot g_i)) \tag{axioms (\ref{eq:pca})}\\
%    &= (\sum_{i=1}^{n+1}p_i\cdot f_i)+_p ( \sum_{i=1}^{n+1}p_i\cdot g_i)\tag{Def. of $\sum_{i=1}^{n+1}{p_i}\cdot f_i$ in (\ref{eq:def somma n pca})}
%\end{align}
        \end{enumerate}
\end{proof}



\section{Appendix to Section \ref{sec:cbproducts}}\label{app:sec:cbproducts}




\subsection{Proofs of Section \ref{ssec:convexproducts}}
\begin{proof}[Proof of Lemma~\ref{lemma:naryconvexproduct}]%\marginpar{Check and move in the Appendix-ok}
By induction on $n$, we show that there exists an object $\bigotimes_{i=1}^n X_i$ which is a convex product of $X_1, \dots, X_n$. For $n=0$, take the final object of $\Cat{C}$ which exists by hypothesis. % that the object $\zero$ is final by Definition~\ref{def:convbicat}.
For $n+1$ assume by induction hypothesis that $\bigotimes_{i=1}^n X_i$ is a convex product of $X_1, \dots, X_n$ with projection $\pi_i\colon \bigotimes_{i=1}^n X_i \to X_i$. Since $\Cat{C}$ has binary convex products we can take the convex product $(\bigotimes_{i=1}^n X_i) \otimes X_{n+1}$ with projections $\pi'_1\colon (\bigotimes_{i=1}^n X_i) \otimes X_{n+1} \to \bigotimes_{i=1}^n X_i$ and $\pi'_2\colon (\bigotimes_{i=1}^n X_i) \otimes X_{n+1} \to X_{n+1}$. We claim that $(\bigotimes_{i=1}^n X_i) \otimes X_{n+1}$ with projections $\pi''_i\colon (\bigotimes_{i=1}^n X_i) \otimes X_{n+1} \to X_i$ defined for all $i= 1,\dots, n+1$ as 
\[\pi''_i = \begin{cases} \pi_1'; \pi_i & i\in 1,\dots, n \\ \pi_2' & i=n+1 \end{cases}\]
is a convex product of $X_1, \dots, X_{n+1}$. 
To check the universal property, let $\vec{p}=p_1,\dots, p_{n+1}$ such that $\sum_{i=1}^{n+1}p_i\leq 1$ and $f_i\colon A \to X_i$. Take $\vec{q}=q_1,\dots, q_n$ with $q_i=\frac{p_i}{1-p_{n+1}}$ if $p_{n+1}\not= 1$ and $q_i=0$ otherwise. Since  $\bigotimes_{i=1}^n X_i$ is an $n$-ary convex product by induction hypothesis, there exists an arrow $\langle f_1, \dots ,f_n\rangle_{\vec{q}} \colon A \to \bigotimes_{i=1}^n X_i$ such that $q_i\cdot f_i = \langle f_1, \dots ,f_n\rangle_{\vec{q}} ; \pi_i$. For the $n+1$-ary, we construct the mediating morphism $h\colon A \to (\bigotimes_{i=1}^n X_i) \otimes X_{n+1}$ as $\langle \, \langle f_1, \dots ,f_n\rangle_{\vec{q}} \,,\, f_{n+1} \, \rangle_{1-p_{n+1},p_{n+1}}$.
 \end{proof}



\begin{proof}[Proof of Proposition \ref{prop:productvsconvex}]%\marginpar{Check and move in the Appendix-ok}
Suppose that $(Z,\pi_1,\pi_2)$ is a convex product of $X_1$ and $X_2$. For all arrows $f_1\colon A \to X_1$ and $f_2\colon A \to X_2$, take the pairing $\langle f_1,f_2\rangle\colon A \to Z$ to be $\langle f_1,f_2\rangle_{p_1,p_2}$ for some $p_1,p_2\in (0,1)$ with $p_1+p_2\leq 1$. Since $p_i\cdot f_i = f_i$ for $i\in\{1,2\}$, then $\langle f_1,f_2\rangle ; \pi_i =  \langle f_1,f_2\rangle_{p_1,p_2};\pi_i= p_i\cdot f_i =f_i$. To show uniqueness, let $h\colon A \to Z$ be such that $h;\pi_i =f_i$. Since $f_i = p_i\cdot f_i $, then $h;\pi_i =p_i\cdot f_i$. By the uniqueness provided by the convex product, $h$ must be $\langle f_1,f_2\rangle_{p_1,p_2}$, that is $\langle f_1,f_2\rangle$. Hence $(Z,\pi_1,\pi_2)$ is a product.

Vice versa, suppose that $(Z,\pi_1,\pi_2)$ is a product of $X_1$ and $X_2$. For all arrows $f_1\colon A \to X_1$ and $f_2\colon A \to X_2$ and $p_1,p_2\in[0,1]$, take $\langle f_1,f_2 \rangle_{p_1,p_2}\colon A \to Z$ to be the pairing $\langle f_1',f_2' \rangle$ where 
\[
f'_i = \begin{cases} 
f_i &\text{if }p_i \neq 0 \\
\star_{A,X_i} &\text{if }p_i=0\end{cases}
\]
Since $p_i\cdot f_i=f_i$ for all $p_i\in (0,1]$ and $p_i \cdot f_i =\star_{A,X_i}$ for $p_i=0$, it holds that for all $p_i\in[0,1]$, $p_i \cdot f_i =f_i'$.
Thus $\langle f_1,f_2 \rangle_{p_1,p_2}; \pi_i = \langle f_1',f_2' \rangle;\pi_i =f_i' =p_i \cdot f_i$. To prove uniqueness, let $h\colon A \to Z$ be such that $h;\pi_i =p_i\cdot f_i$. Since $p_i\cdot f_i = f_i'$, then $h;\pi_i = f_i '$. 
By the uniqueness provided by the product, $h$ must be $\langle f_1',f_2'\rangle$, that is $\langle f_1,f_2 \rangle_{p_1,p_2}$. Hence, $(Z,\pi_1,\pi_2)$ is a convex product.
%\[\langle f_1,f_2 \rangle_{p_1,p_2} = \begin{cases} 
%\langle f_1,f_2\rangle & \text{if} p_1\neq 0, \; p_2\neq 0\\
%\langle \star_{A,X_1},f_2\rangle & \text{if} p_1= 0, \; p_2\neq 0\\
%\langle f_1, \star_{A,X_2}\rangle & \text{if} p_1\neq 0, \; p_2= 0\\
%\langle \star_{A,X_1}, \star_{A,X_2}\rangle & \text{if} p_1\neq 0, \; p_2= 0\\
%\end{cases}\]
\end{proof}

%
%
%
%
%
%

\subsection{Proof of Section \ref{ssec:convexbiproduct}}

\begin{proof}[Proof of Lemma~\ref{lemma:initialproperties}]
We prove below item by item.
\begin{enumerate}
\item Follows easily by enrichment and finality of $\zero$: 
\begin{align*}
\star_{X,Y} & = \star_{X,\zero} ; \cobang{Y} \tag{\ref{eq:enr}}\\
&= \bang{X};\cobang{Y} \tag{$\zero$ is final}
\end{align*}
\item Consider the arrows $f;\iota_1$ and $g;\iota_2$ from $X$ to $Y\oplus Y$ and their sum $f;\iota_1 +_p g;\iota_2$. The following computation shows that 
    \begin{align}
        (f;\iota_1 +_p g;\iota_2);\pi_1 & = f;\iota_1;\pi_1 +_p g;\iota_2;\pi_1 \tag{\ref{eq:enr}}\\
        & = f +_p \star_{X,Y} \tag{by \eqref{eq:delta}}\\
        & = p \cdot f \tag{def. of $p\cdot -$}
    \end{align} 
    and similarly $(f;\iota_1 +_p g;\iota_2);\pi_2 = (1-p)\cdot g$. By the universal property of convex products, we conclude that $f;\iota_1 +_p g;\iota_2 = \langle f,g \rangle_{p,1-p}$. Postcomposing both sides with $[\id{Y},\id{Y}]$ we obtain 
    \begin{align}
        \langle f,g \rangle_{p,1-p};[\id{Y},\id{Y}] & = (f;\iota_1 +_p g;\iota_2);[\id{Y},\id{Y}] \notag\\
        &= (f;\iota_1 ;[\id{Y},\id{Y}] +_p g;\iota_2 ;[\id{Y},\id{Y}]) \tag{\ref{eq:enr}}\\
        &=(f ;\id{Y} +_p g ; \id{Y}) \tag{properties of $[-,-]$}\\
        & = f +_p g.   \tag{category}
    \end{align}
\item We prove that $ \id{X_1}\oplus \cobang{X_2}=\runit{X_1}; \iota_1$ by universal property of $X_1\oplus \zero$ using the injections $\iota_1^{X_1\oplus\zero} \colon X_1 \to X_1\oplus \zero$ and 
$\iota_2^{X_1\oplus\zero} \colon \zero \to X_1\oplus \zero$. First observe that, by initiality of $\zero$,
 \[\iota_2^{X_1\oplus \zero}; (\id{X_1}\oplus \cobang{X_2}) = \iota_2^{X_1\oplus \zero}; \runit{X_1}; \iota_1 \text{.} \]
Moreover
\begin{align*}
\iota_1^{X_1\oplus \zero} ; (\id{X_1}\oplus \cobang{X_2}) &= \iota_1^{X_1\oplus \zero} ; [\id{X_1}; \iota_{1},\cobang{X_2};\iota_2 ] \tag{def. of $\oplus$}\\
&=\id{X_1};\iota_1 \tag{coproduct}\\
&=\iota_1^{X_1\piu\zero}; \runit{X_1} ; \iota_1. \tag{def. of $\runit{}$} 
\end{align*}
Thus $ \id{X_1}\oplus \cobang{X_2}=\runit{X_1}; \iota_1$.
\item Analogous to the point above.
\item Observe that, for $\iota_i\colon X_i \to X_1\oplus X_2$, 
\[\iota_1 ; (\id{X_1}\oplus \bang{X_2});\runit{X_1}=\id{X_1} \qquad \iota_2 ;  (\id{X_1}\oplus \bang{X_2});\runit{X_1} = \bang{X_2}; \cobang{X_1}= \star_{X_2,X_1}\]
namely $(\id{X_1}\oplus \bang{X_2});\runit{X_1} = [\id{X_1}, \star_{X_2,X_1}]$. Now, by \eqref{eq:delta}, one has that $[\id{X_1}, \star_{X_2,X_1}]=\pi_1$.
\item Analogous to the point above.
\item By the following derivation.
\begin{align*}
p \cdot f &= f+_p\star_{X,Y} \tag{def. of $p\cdot(-)$}\\
&=\, \diagp{X} ; (f \oplus \star_{X,Y}) ; \codiag{Y} \tag{Lemma~\ref{lemma:initialproperties}.\ref{eq: assioma aggiuntivo}} \\ 
&=\, \diagp{X} ; (f \oplus (\bang{X};\cobang{Y})) ; \codiag{Y} \tag{Lemma \ref{lemma:initialproperties}.\ref{lemma:starxy}} \\ 
&=\, \diagp{X} ; (f \oplus \bang{X}) ;(\id{Y} \oplus \cobang{Y}) ; \codiag{Y} \tag{Symmetric Monoidal Category} \\ 
&=\, \diagp{X} ; (f \oplus \bang{X}) ;\runit{Y}. \tag{\ref{ax:Mon2}} \\ 
 \end{align*}
 \item By the following derivation.
\begin{align*}
(1-p) \cdot f &= f+_{1-p}\star_{X,Y} \tag{def. of $p\cdot(-)$}\\
&= \star_{X,Y} +_{p} f \tag{\ref{eq:pca}}\\
&=\, \diagp{X} ; (\star_{X,Y} \oplus f) ; \codiag{Y} \tag{Lemma~\ref{lemma:initialproperties}.\ref{eq: assioma aggiuntivo}} \\ 
&=\, \diagp{X} ; ((\bang{X};\cobang{Y}) \oplus f) ; \codiag{Y} \tag{Lemma \ref{lemma:initialproperties}.\ref{lemma:starxy}} \\ 
&=\, \diagp{X} ; (\bang{X} \oplus f) ;(\id{Y} \oplus \cobang{Y}) ; \codiag{Y} \tag{Symmetric Monoidal Category} \\ 
&=\, \diagp{X} ; ( \bang{X} \oplus f) ;\lunit{Y}. \tag{\ref{ax:Mon2}} \\ 
 \end{align*} 
 
\end{enumerate}
\end{proof}





\begin{lemma}\label{lemma:naturality}
In a convex biproduct category, $\bang{}$ and $\diagp{}$ form natural transformations, namely 
$f;\bang{Y}=\bang{X}$  and $f;\diagp{Y} =\, \diagp{X};(f\oplus f)$ for all $f\colon X\to Y$.
\end{lemma}
\begin{proof}
The equation $f;\bang{Y}=\bang{X}$ for $f\colon X \to Y$ follows from the fact that $\zero$ is final. The equation
\[f;\diagp{Y} =\, \diagp{X};(f\oplus f)\]
follows from the universal property of convex products. Indeed, observe that
\begin{align*}
    \diagp{X};(f\oplus f);\pi_1 &=\, \diagp{X};(f\oplus f);(\id{Y}\piu \bang{Y});\runit{X} \tag{Lemma \ref{lemma:initialproperties}.\ref{lemma:pi1}}\\
    &=\, \diagp{X};(f\oplus (f;\bang{Y}));\runit{X}  \tag{Symmetric Monoidal Category}\\
    &=\, \diagp{X};(f\oplus \bang{X});\runit{X} \tag{Naturality of $\bang{X}$}\\
    &= p\cdot f \tag{Lemma \ref{lemma:initialproperties}.\ref{lemma:pf}}
\end{align*}
\begin{align*}
    \diagp{X};(f\oplus f);\pi_2 &=\, \diagp{X};(f\oplus f);(\bang{Y} \piu \id{Y} );\lunit{X} \tag{Lemma \ref{lemma:initialproperties}.\ref{lemma:pi2}}\\
    &=\, \diagp{X};( (f;\bang{Y})\oplus f);\lunit{X}  \tag{Symmetric Monoidal Category}\\
    &=\, \diagp{X};( \bang{X} \oplus f);\lunit{X} \tag{Naturality of $\bang{X}$}\\
    &= (1-p)\cdot f \tag{Lemma \ref{lemma:initialproperties}.\ref{lemma:1-pf}}
\end{align*}
and that
\begin{align*}
f;\diagp{Y};\pi_1 &= f ; \langle \id{X}, \id{X} \rangle_{p,(1-p)};\pi_1 \tag{def. of $\diagp{}$}\\
&= f; (p\cdot \id{X}) \tag{convex product} \\
&= p\cdot(f;\id{X}) \tag{Lemma \ref{lemma:initialproperties2}.\ref{lemma:p;}}\\
&=p\cdot f \tag{Category}
\end{align*}
\begin{align*}
f;\diagp{Y};\pi_2 &= f ; \langle \id{X}, \id{X} \rangle_{p,(1-p)};\pi_2 \tag{def. of $\diagp{}$}\\
&= f; (\,(1-p)\cdot \id{X}\,) \tag{convex product} \\
&= (1-p)\cdot(f;\id{X}) \tag{Lemma \ref{lemma:initialproperties2}.\ref{lemma:p;}}\\
&=(1-p)\cdot f \tag{Category}
\end{align*}
\end{proof}

\begin{lemma}\label{lemma:coherence}
In a convex biproduct category, the coherence axioms in Figure~\ref{fig:copcacoherence} hold.
\end{lemma}
\begin{proof}
(Coh2), (Coh3) and (Coh4) follow from the fact that $\zero$ is both initial and final. (Coh1) follows from the naturality and the universal property of coproducts. 

Let $s\defeq \assoc{X}{X}{Y\piu Y}; ( {\id{X}\piu \Iassoc{X}{Y}{Y}} ); ({\id X \piu ( \symm{X}{Y}^{\piu} \piu \id Y)} ); (\id{X}\piu \assoc{Y}{X}{Y}); \Iassoc{X}{Y}{X\oplus Y}$. Let
\[\iota_1^{(X\oplus X)\oplus(Y\oplus Y)}\colon X\piu X \to (X\piu X)\piu (Y\piu Y) \text{ and }\iota_2^{(X\oplus X)\oplus(Y\oplus Y)}\colon Y\piu Y \to (X\piu X)\piu (Y\piu Y)\] be the injections. By using the definitions of the structural isomorphisms induced by the coproducts, one can easily check that
\begin{equation}\label{eq:local}
\iota_1^{(X\oplus X)\oplus(Y\oplus Y)}; s= \iota_1 \oplus \iota_1 \qquad \text{and} \qquad \iota_2^{(X\oplus X)\oplus(Y\oplus Y)}; s= \iota_2 \oplus \iota_2\text{.}
\end{equation}
Then
\begin{align*}
\iota_1; (\diagp{X}\piu\, \diagp{Y}) ;s &= \iota_1; [\diagp{X}; \iota_1^{(X\oplus X)\oplus(Y\oplus Y)} \, , \, \diagp{Y}; \iota_2^{(X\oplus X)\oplus(Y\oplus Y)}] ;s \tag{def. of $\oplus$}\\
& =\, \diagp{X};\iota_1^{(X\oplus X)\oplus(Y\oplus Y)};s  \tag{coproduct}\\
&=\, \diagp{X}; (\iota_1 \oplus \iota_1) \tag{\ref{eq:local}} \\
&= \iota_1;\,\diagp{X\piu Y} \tag{Lemma \ref{lemma:naturality}}
\end{align*}
and 
\begin{align*}
\iota_2; (\diagp{X}\piu\, \diagp{Y}) ;s &= \iota_2; [\diagp{X}; \iota_1^{(X\oplus X)\oplus(Y\oplus Y)} \, , \, \diagp{Y}; \iota_2^{(X\oplus X)\oplus(Y\oplus Y)}] ;s \tag{def. of $\oplus$}\\
& =\, \diagp{Y};\iota_2^{(X\oplus X)\oplus(Y\oplus Y)};s  \tag{coproduct}\\
&=\, \diagp{Y}; (\iota_2 \oplus \iota_2) \tag{\ref{eq:local}} \\
&= \iota_2;\,\diagp{X\piu Y} \tag{Lemma \ref{lemma:naturality}}
\end{align*}
Hence, the universal property of coproducts implies that $\diagp{X\piu Y} = \diagp{X}\piu \diagp{Y};s$.
\end{proof}


\begin{lemma}\label{lemma:pdiag} \label{lemma:diagp}
In a convex biproduct category the following hold:
\begin{enumerate}
\item $p \cdot \diagq{X} = \langle \id{X},\id{X}\rangle_{pq, p(1-q)} $.%\label{lemma:pdiag}
\item $\diagp{X}; (\id{X} \oplus q\cdot \id{X})=\langle \id{X},\id{X}\rangle_{p,(1-p)q}$ %\label{lemma:diagp}
\end{enumerate}
\end{lemma}
\begin{proof}
We prove below item by item.
\begin{enumerate}
\item By the uniqueness of $\langle \id{X},\id{X}\rangle_{pq, p(1-q)}$ and the following two derivations.
\begin{align*}
(p \cdot \diagq{X})  ; \pi_1 &= p \cdot (\diagq{X};\pi_1) \tag{Lemma \ref{lemma:initialproperties}.4}\\
&=p\cdot (\langle \id{X},\id{X}\rangle_{q,1-q} ; \pi_1) \tag{def. of $\diagq{}$}\\
&= p\cdot (q \cdot \id{X}) \tag{convex product}\\
&= (p q) \cdot \id{X} \tag{Lemma~\ref{lemma:initialproperties2}.\ref{lemma:pq}}
\end{align*}
\begin{align*}
(p \cdot \diagq{X})  ; \pi_2 &= p \cdot (\diagq{X};\pi_2) \tag{Lemma \ref{lemma:initialproperties}.4}\\
&=p\cdot (\langle \id{X},\id{X}\rangle_{q,1-q} ; \pi_2) \tag{def. of $\diagq{}$}\\
&= p\cdot ((1-q) \cdot \id{X}) \tag{convex product}\\
&= (p (1-q)) \cdot \id{X} \tag{Lemma~\ref{lemma:initialproperties2}.\ref{lemma:pq}}
\end{align*}
\item By uniqueness of $\langle \id{X},\id{X}\rangle_{p,(1-p)q}$ and the following two derivations. 
\begin{align*}
\diagp{X}; (\id{X} \oplus q\cdot \id{X}) ; \pi_1 &=\, \diagp{X}; (\id{X} \oplus q\cdot \id{X}) ; (\id{X}\oplus \bang{X});\runit{X} \tag{Lemma \ref{lemma:initialproperties}.\ref{lemma:pi1}}\\
&=\, \diagp{X}; (\id{X} \oplus (q\cdot \id{X} ; \bang{X})); \runit{X} \tag{Symmetric Monoidal Category}\\
&=\, \diagp{X}; (\id{X} \oplus  \bang{X}); \runit{X} \tag{Naturality of $\bang{}$} \\
&= p\cdot \id{X} \tag{Lemma \ref{lemma:initialproperties}.\ref{lemma:pf}}
\end{align*}
\begin{align*}
\diagp{X}; (\id{X} \oplus q\cdot \id{X}) ; \pi_2 &=\, \diagp{X}; (\id{X} \oplus q\cdot \id{X}) ; (\bang{X}\oplus \id{X});\lunit{X} \tag{Lemma \ref{lemma:initialproperties}.\ref{lemma:pi2}}\\
&=\, \diagp{X}; ( \bang{X}\oplus( q\cdot \id{X}) ) ; \lunit{X} \tag{Symmetric Monoidal Category}\\
&= (1-p)\cdot (q \cdot \id{X}) \tag{Lemma \ref{lemma:initialproperties}.\ref{lemma:1-pf}}\\
&= ((1-p)q) \cdot \id{X}. \tag{Lemma \ref{lemma:initialproperties2}.\ref{lemma:pq}}
\end{align*}
\end{enumerate}
\end{proof}




\begin{proof}[Proof of Proposition \ref{lemma: copca objects in convbicat}]
The first point follows immediately by the Fox theorem \cite{fox1976coalgebras}.

For the second point, we proved naturality and coherence in Lemmas~\ref{lemma:naturality} and \ref{lemma:coherence}. Below we prove the axioms in Figure~\ref{fig:co-pca axioms}.

    (PCA1) in Figure~\ref{fig:co-pca axioms} follows from the universal property of convex products. Indeed, observe that 
    \begin{align*}
        \diagp{X} ; (\diagq{X} \oplus \id{X}) ; \pi_1 &=\, \diagp{X} ; (\diagq{X} \oplus \id{X}) ;( \id{X \oplus X} \oplus \bang{X}) ; \runit{X\oplus X} \tag{Lemma \ref{lemma:initialproperties}.\ref{lemma:pi1}}\\
        &=\, \diagp{X} ; (\diagq{X} \oplus \bang{X}) ; \runit{X\oplus X} \tag{Symmetric Monoidal Category}\\
        &= p \cdot \diagq{} \tag{Lemma \ref{lemma:initialproperties}.\ref{lemma:pf}}\\
        &= \langle \id{X}, \id{X} \rangle_{pq,\,p(1-q)}\tag{Lemma \ref{lemma:pdiag}}
    \end{align*}
    
        \begin{align*}
        \diagp{X} ; (\diagq{X} \oplus \id{X}) ; \pi_2 &=\, \diagp{X} ; (\diagq{X} \oplus \id{X}) ;( \bang{X \oplus X} \oplus \id{X}) ; \lunit{X} \tag{Lemma \ref{lemma:initialproperties}.\ref{lemma:pi2}}\\
        & =\, \diagp{X} ; ( \, (\diagq{X} ; (\bang{X} \oplus \bang{X}) ) \oplus \id{X} \,)  ; \lunit{X} \tag{Symmetric Monoidal Category}\\
        &=\, \diagp{X} ;  (\bang{X}  \oplus \id{X} )  ; \lunit{X} \tag{Lemma \ref{lemma:naturality}} \\
        &=  (1-p)\cdot \id{X} \tag{Lemma \ref{lemma:initialproperties}.\ref{lemma:1-pf}}
    \end{align*}
    
Similarly, by fixing  $\tilde{p}\defeq  pq$ and $\tilde{q} \defeq \frac{p(1-q)}{1-pq}$, it holds  that
\begin{align}
    \diagptilde{}; (\id{X}\oplus\, \diagqtilde{}) ; \Iassoc{X}{X}{X}; \pi_1 &=\, \diagptilde{}; (\id{X}\oplus\,  \diagqtilde{}) ; \Iassoc{X}{X}{X}; ( \id{X \oplus X} \oplus \bang{X}) ; \runit{X\oplus X} \tag{Lemma \ref{lemma:initialproperties}.\ref{lemma:pi1}}\\
     &=\, \diagptilde{}; (\id{X}\oplus\, \diagqtilde{}) ; ( \id{X} \oplus (\id{X} \oplus \bang{X})) ; \runit{X} \tag{Symmetric Monoidal Category}\\
 &=\, \diagptilde{}; (\,\id{X}\oplus ( \diagqtilde{} ;(\id{X} \oplus \bang{X})) \,); \runit{X} \tag{Symmetric Monoidal Category}\\
&=\, \diagptilde{}; (\,\id{X}\oplus \tilde{q}\cdot \id{X} \,) \tag{Lemma \ref{lemma:initialproperties}.\ref{lemma:pf}}\\
    &=\langle \id{X}, \id{X} \rangle_{\tilde{p},(1-\tilde{p})\tilde{q}} \tag{Lemma \ref{lemma:diagp}}\\
    &= \langle \id{X}, \id{X} \rangle_{pq,p(1-q)} \tag{def.\ of $\tilde{p}$ and $\tilde{q}$}
\end{align} 
    \begin{align}%\marginpar{Da finire qua, mi sembra che si debba usare la coerenza}
    \diagptilde{}; (\id{X}\oplus\, \diagqtilde{}); \Iassoc{X}{X}{X};\pi_2&=\, \diagptilde{};( \id{X}\oplus\, \diagqtilde{}); \Iassoc{X}{X}{X}; (\bang{X\piu X}\piu \id{X});\lunit{X}  \tag{Lemma \ref{lemma:initialproperties}.\ref{lemma:pi2}}\\
    &=\, \diagptilde{};( \id{X}\oplus\, \diagqtilde{}); \Iassoc{X}{X}{X}; ((\bang{X}\piu \bang{X})\piu \id{X});\lunit{X}  \tag{Lemma \ref{lemma:coherence}}\\
    &=\, \diagptilde{}; (\bang{X} \oplus \id{X}); \lunit{X} ; \diagqtilde{X}; (\bang{X}\piu \id{X}) ;\lunit{X}  \tag{Symmetric Monoidal Category}\\
    &=\, \diagptilde{}; (\bang{X} \oplus \id{X}); \lunit{X} ; (1-\tilde{q})\cdot \id{X}  \tag{Lemma \ref{lemma:initialproperties}.\ref{lemma:1-pf}}\\
    &=\, \diagptilde{}; (\bang{X} \oplus ( (1-\tilde{q})\cdot \id{X} )) ;  \lunit{X} \tag{Symmetric Monoidal Category}\\
    &= (1-\tilde{p}) \cdot ((1-\tilde{q})\cdot \id{X} )  \tag{Lemma \ref{lemma:initialproperties}.\ref{lemma:1-pf}}\\
    &=  (1-\tilde{p}) (1-\tilde{q}) \cdot \id{X} \tag{Lemma \ref{lemma:initialproperties2}.\ref{lemma:pq}}\\
    &= (1-p) \cdot \id{X} \tag{def. $\tilde{p}$ and $\tilde{q}$}
\end{align}
    Hence, $\diagp{} ; (\diagq{} \oplus \id{X}) =\, \diagptilde{} ; \id{X}\oplus\, \diagqtilde{}; \Iassoc{X}{X}{X}$.\\

    (PCA2) in Figure~\ref{fig:co-pca axioms} follows from the coherence of the PCA-enrichment (i.e.,  Lemma~\ref{lemma:initialproperties}.\ref{eq: assioma aggiuntivo}):
  \begin{align*}
    \diagp{X}; \codiag{X} &=\,   \diagp{X}; (\id{X} \oplus \id{X}) \codiag{X}\\
    &= \id{X}+_p \id{X} \tag{Lemma~\ref{lemma:initialproperties}.\ref{eq: assioma aggiuntivo}}\\
    &=\id{X} \tag{\ref{eq:pca}}
    \end{align*}
    \\

    (PCA3) in Figure~\ref{fig:co-pca axioms} follows from the universal property of convex products. Indeed, observe that
    \begin{align*}
        \diagp{X};\symm{X}{X};\pi_1 &=\, \diagp{X};\symm{X}{X};(\id{X}\piu \bang{X});\runit{X} \tag{Lemma \ref{lemma:initialproperties}.\ref{lemma:pi1}}\\
        &=\, \diagp{X};(\bang{X} \piu \id{X});\lunit{X} \tag{Symmetric Monoidal Category}\\
        &= (1-p)\cdot \id{X}. \tag{Lemma \ref{lemma:initialproperties}.\ref{lemma:1-pf}}
    \end{align*}
    And similarly for $\pi_2$ we have that $\diagp{X};\symm{X}{X};\pi_2 = p\cdot \id{X}$. Hence, the universal property of convex products implies that $\diagp{X};\symm{X}{X} =\, \diagpbar{X}$.
 \end{proof}

%
%
%
%
%

\begin{proposition}\label{prop: functor1}
    Let $F\colon \Cat{C}\to \Cat{D}$ be a morphism between convex biproduct categories. 
     $F$ preserves convex products.
   \end{proposition}
   \begin{proof}[Proof of Proposition \ref{prop: functor1}]
    Preservation of finite coproducts is equivalent to preservation of binary coproducts and initial object. Hence, assume that $F(X)\overset{F(\iota_1)}{\to}F(X\piu Y)\overset{F(\iota_2)}{\leftarrow}$ is a coproduct of $F(X)$ and $F(Y)$, which is equivalent to requiring that the arrow $[F(\iota_1),F(\iota_2)]\colon F(X)\piu F(Y)\to F(X\piu Y)$ is an isomorphism (call $k$ the inverse). And assume also that $F(\zero_\Cat{C})$ is initial in $\Cat{D}$ (hence $\cobang{F(\zero_{\Cat{C}})}$ is an isomorphism). Now observe that
    \begin{align}
        F(\pi_1^{X\piu Y})&= F(\id{X}\piu\, \bangp{Y};\rho_X) \tag{Lemma \ref{lemma:initialproperties}.\ref{lemma:pi1}}\\
        &=F(\id{X}\piu\, \bangp{Y});F([\id{X}, \cobang{X}]) \tag{def. $\rho$}\\
        &= F([\id{X}, \bangp{Y};\cobang{X}]) \tag{Coproducts}\\
        &=k;[\id{F(X)}, F(\,\bangp{Y};\cobang{X})] \tag{coproducts and $k$ iso}\\
        &= k;(\id{F(X)}\oplus\, \bangp{F(Y)});[\id{F(X)},\cobang{F(X)}] \tag{$F(\zero_\Cat{C})\cong\zero_\Cat{D}$}\\
        &= k;\pi_1^{F(X)\piu F(Y)} \tag{def.\ $\rho$}
    \end{align}
    and similarly $F(\pi_2^{X\piu Y})=k;\pi_2^{F(X)\piu F(Y)}$. We can now prove that $F(X\piu Y)$ with projections $F(\pi_1^{X\piu Y})$ and $F(\pi_2^{X\piu Y})$ is a convex binary product of $F(X)$ and $F(Y)$. Let $f\colon A\to F(X)$ and $g\colon A\to F(Y)$ and $p\in (0,1)$ and take the arrow $\langle f,g \rangle_{p,1-p};[F(\iota_1), F(\iota_2)]\colon A \to F(X \piu Y)$, where $\langle f,g \rangle_{p,1-p}\colon A \to F(X)\piu F(Y)$ is the unique arrow induced by convex binary products in $\Cat{D}$. Now, since $F(\pi_1^{X\piu Y})= k; \pi_1^{F(X\piu F(Y))}$ it follows that 
    \begin{align}
        \langle f,g \rangle_{p,1-p};[F(\iota_1), F(\iota_2)];F(\pi_1^{X\piu Y})&= \langle f,g \rangle_{p,1-p};[F(\iota_1), F(\iota_2)];k;\pi_1^{F(X)\piu F(Y)}  \tag*{}\\
        &= \langle f,g \rangle_{p,1-p};\pi_1^{F(X)\piu F(Y)} \tag{$k$ inverse}\\
        &=p\cdot f \tag{convex products}
    \end{align}
    and similarly $\langle f,g \rangle_{p,1-p};[F(\iota_1), F(\iota_2)];F(\pi_2^{X\piu Y})=(1-p)\cdot g$. For the uniqueness, if $h:A \to F(X\piu Y)$ is such that $h;F(\pi_1^{X\piu Y})=p\cdot f$ and $h;F(\pi_2^{X\piu Y})=(1-p)\cdot g$ then $h;k= \langle f,g \rangle_{p,1-p}$. Post-composition with $[F(\iota_1),F(\iota_2)]$ gives $h= \langle f,g \rangle_{p,1-p};[F(\iota_1), F(\iota_2)]$.
   \end{proof}
   
   \begin{proof}[Proof of Proposition \ref{prop:functor 1 e mezzo}]
    It follows from Proposition~\ref{prop: functor1} and Lemma~\ref{lemma:naryconvexproduct} which shows how to produce n-ary convex products from binary ones.
\end{proof}
   

\begin{proof}[Proof of Proposition \ref{prop: monoidal functors}]\label{proof: prop monoidal functors}
By Fox's theorem we know that every finite coproduct category $\Cat{C}$ (hence any convex biproduct category) correspond to  monoidal category $(\Cat{C},\oplus_\Cat{C},\zero_\Cat{C})$ in which every object $X$ is equipped with a coherent and natural monoid structure $(X,\codiag{X},\cobang{X})$. Moreover, a functor $F\colon \Cat{C}\to \Cat{D}$ between finite coproduct  categories preserves finite coproducts if and only if seen as a monoidal functor $F\colon(\Cat{C},\oplus_\Cat{C},\zero_\Cat{C})\to (\Cat{D},\oplus_\Cat{D},\zero_\Cat{D})$ it is strong monoidal  and preserves monoids
    \begin{equation*}
        \psi_{X,X};F(\codiag{X})= \codiag{F(X)} \qquad \psi_\zero;F(\cobang{X})=\cobang{F(X)}
    \end{equation*}
    where  $\psi\colon  \piu_\Cat{C}\circ (F\times F) \overset{\cong}{\Rightarrow} F\circ \piu_{\Cat{D}}$ and $\psi_\zero\colon \zero_\Cat{D}\overset{\cong}{\to} F(\zero_\Cat{C})$ are given by the strong monoidal structure. For any convex biproduct category,  Proposition~\ref{lemma: copca objects in convbicat} implies that every object $X$ is equipped with a natural and coherent co-pca structure $(X,\diagp{X}, \bangp{X})$. In this case, a strong monoidal functor that preserves monoids between monoidal categories, in which every object has the structure of a coherent and natural monoid and co-pca, also preserves $\,\bangp{X}$. Indeed,
     \begin{align}
        \psi_0^{-1}&= \psi_0^{-1};\,\bangp{0} \tag{\ref{ax:coh3} in Figure~\ref{fig:copcacoherence} }\\
        &=\bangp{F(0)} \tag{\ref{eq:nat copca1}}
    \end{align}
    Hence, $\,\bangp{X}$ is preserved: $F(\,\bangp{X});\psi_\zero^{-1}=F(\,\bangp{X});\bangp{F(0)}=\bangp{F(X)}$. 
    
    Moreover, observe that a convex biproduct category $\Cat{C}$ corresponds to a monoidal category 
    $(\Cat{C},\oplus_\Cat{C},\zero_\Cat{C})$ in which every object is equipped with coherent and natural monoid and co-pca structures and moreover it holds that given $f:X\to A$ and $g:X\to B$ then for every $h\colon X\to A\piu B$
    \begin{equation}\label{eq:ax aggiuntivo per fox convex biprod}
         \begin{cases}
    h;(\id{A}\piu\, \bangp{B})=\, \diagp{X};(f\piu \,\bangp{B})\\
    h;(\,\bangp{A}\piu \id{B})=\,\diagp{X};(\,\bangp{A}\piu g)
\end{cases}
\qquad
\Rightarrow 
\qquad 
h=\,\diagp{X};(f\piu g)
    \end{equation}
    since $\Cat{C}$ has convex binary products.
    Now if $F$ is PCA-enriched, then 
    \begin{align}
        F(\diagp{X});\psi_{X,X}^{-1};(\id{F(X)}\piu\, \bangp{F(X)})&=F(\diagp{X});\psi_{X,X}^{-1};(\id{F(X)}\piu\, \bangp{F(X)}); \rho_{F(X)};\rho^{-1}_{F(X)} \tag*{}\\
        &=F(\diagp{X});\psi_{X,X}^{-1};(\id{F(X)}\piu\, \bangp{F(X)}); (\id{F(X)}\piu \cobang{F(X)});\codiag{F(X)};\rho_{F(X)}^{-1} \tag{Def. $\rho$}\\
        &=F(\diagp{X});F(\id{X}\piu (\,\bangp{X};\cobang{X})); \psi_{X,X}^{-1};\codiag{F(X)};\rho^{-1}_{F(X)}\tag{Nat. $\psi$}\\
        &=F(\diagp{X});F(\id{X}\piu (\,\bangp{X};\cobang{X}));F(\codiag{X});\rho^{-1}_{F(X)} \tag{$\psi_{X,X};F(\codiag{X})= \codiag{F(X)} $}\\
        &= F(p\cdot \id{X});\rho^{-1}_{F(X)}  \tag{Lemma \ref{lemma:initialproperties}.\ref{eq: assioma aggiuntivo}}\\
        &= p\cdot \id{F(X)};\rho^{-1}_{F(X)}\tag{PCA-enrichment}\\
        &=\, \diagp{F(X)};(\id{F(X)}\piu\,\bangp{F(X)});\rho_{F(X)};\rho^{-1}_{F(X)}\tag{Lemma \ref{lemma:initialproperties}.\ref{eq: assioma aggiuntivo} + Def. $\rho$ + Lemma \ref{lemma:initialproperties}.\ref{lemma:starxy}}\\
        &=\, \diagp{F(X)};(\id{F(X)}\piu\,\bangp{F(X)}) \tag*{}
    \end{align}
and similarly one can prove that $F(\diagp{X});\psi_{X,X}^{-1};(\,\bangp{F(X)}\piu\id{F(X)})= \, \diagp{F(X)};(\,\bangp{F(X)}\piu \id{F(X)})$. Hence, by property (\ref{eq:ax aggiuntivo per fox convex biprod}) it follows that $F(\diagp{X});\psi_{X,X}^{-1}=\, \diagp{F(X)}$.
Vice versa, if $F\colon(\Cat{C},\oplus_\Cat{C},\zero_\Cat{C})\to (\Cat{D},\oplus_\Cat{D},\zero_\Cat{D})$ also preserves diagonals $\diagp{}$ then it is PCA-enriched since
\begin{align}
    F(f+_pg)&= F(\diagp{X};(f\piu g);\codiag{X}) \tag{Lemma \ref{lemma:initialproperties}.\ref{eq: assioma aggiuntivo}}\\
    &=\, \diagp{F(X)};\psi_{X,X};F(f\piu g; \codiag{X}) \tag*{}\\
&=\, \diagp{F(X)};F(f)\piu F(g);\codiag{F(X)}\tag*{}\\
&= F(f)+_p F(g) \tag{Lemma \ref{lemma:initialproperties}.\ref{eq: assioma aggiuntivo}}
\end{align}
 Similarly, the preservation of $\,\bangp{X}$ and $\cobang{Y}$ implies that $ F(\star_{X,Y})= \star_{F(X),F(Y)}$.

\end{proof}














\begin{comment}
    

\marginpar{La prop \ref{prop: functor2} non la usiamo mai.}
   \begin{proposition}\label{prop: functor2}
     Let $F\colon \Cat{C}\to \Cat{D}$ a functor between convex biproduct categories that  preserves finite coproducts, then $F$ is PCA-enriched if and only if $F$ preserves $\diagp{X}$ for every $X\in\Cat{C}$.
   \end{proposition}
   \begin{proof}
    Using the notation of proposition~\ref{prop: functor1}, if $F$ is PCA-enriched, then 
    \begin{align}
        F(\langle \id{X},\id{X}\rangle_{p,1-p});k;\pi_1^{F(X)\piu F(Y)}&= F(\langle \id{X},\id{X}\rangle_{p,1-p});F(\pi_1^{X\piu Y}) \tag{Prop. \ref{prop: functor1}}\\
        &= F(p\cdot \id{X})= p\cdot \id{F(X)}\tag{PCA-enrichment}
    \end{align}
and similarly $F(\langle id,id\rangle_{p,1-p});k;\pi_1^{F(X)\piu F(Y)}= (1-p)\cdot \id{F(X)}$. Hence, the universal property of convex products in $\Cat{D}$ implies that $F(\langle \id{X},\id{X}\rangle_{p,1-p});k= \langle \id{F(X)}, \id{F(X)} \rangle_{p,1-p}$.

Vice versa, if $F(\langle \id{X},\id{X}\rangle_{p,1-p});k= \langle \id{F(X)}, \id{F(X)} \rangle_{p,1-p}$ then for every pair of arrows $f,g\in \Cat{C}([X,Y])$
\begin{align}
    F( f+_p g)=& F(\langle \id{X},\id{X} \rangle_{p,1-p};(f\piu g);\codiag{Y}) \tag{Lemma \ref{lemma:initialproperties}.\ref{eq: assioma aggiuntivo}}\\
    &= F(\langle \id{X},\id{X} \rangle_{p,1-p});k; (F(f)\piu F(g)); \codiag{F(Y)} \tag{coproducts preserv.}\\
    &=\langle \id{F(X)}, \id{F(X)} \rangle_{p,1-p};(F(f)\piu F(g)); \codiag{F(Y)} \tag*{}\\
    &= F(f)+_pF(g). \tag{Lemma \ref{lemma:initialproperties}.\ref{eq: assioma aggiuntivo}}
\end{align}
Moreover,
\begin{align}
    F(\star_{X,Y})&= F(\,\bangp{X};\cobang{Y}) \tag{Lemma \ref{lemma:initialproperties}.\ref{lemma:starxy}}\\
    &= F(\,\bangp{X});F(\,\cobang{Y}) \tag*{}\\
    &= \bangp{F(X)};\cobang{F(Y)} \tag{$\zero_\Cat{D}$ terminal + $F(\zero_\Cat{C})\cong \zero_\Cat{D}$   }\\
    &= \star_{F(X),F(Y)}. \tag{Lemma \ref{lemma:initialproperties}.\ref{lemma:starxy}}
\end{align}
   \end{proof}

\end{comment}



   
   
   \subsection{Proofs of Section \ref{ssec:almost}}
   
   \begin{proof}[Proof of Lemma \ref{lemma:enrichmentpca}]
We first prove that for all objects $X,Y$ in $\Cat{C}$, the structure defined above is a pca, namely that the equalities in \eqref{eq:pca} hold.
   For idempotency, we have the following derivation. % follows from the following equalities:
   \begin{align}
        f +_p f & =\, \diagp{X};(f \oplus f);\codiag{Y} \tag{def. $+_p$}\\
        & =\, \diagp{Y};\codiag{Y};f \tag{\ref{eq:nat monoid1}}\\
        & = f. \tag{\ref{ax:PCA2}}
   \end{align} 
For commutativity, we have the following derivation.
    \begin{align}
        f +_p g & =\, \diagp{X};(f \oplus g);\codiag{Y} \tag{def. $+_p$}\\
        & =\, \diagp{X};\sigma_{X,X};(g \oplus f);\codiag{Y} \tag{SMC}\\
        & =\, \diaggen{1-p}{X} ;(g \oplus f);\codiag{Y} \tag{\ref{ax:PCA3}}\\
        & = g +_{1-p} f. \tag{def. $+_{1-p}$}
    \end{align}
    For associativity, consider $f,g,h\colon X \to Y$ and $p,q\in[0,1]$. We have:
    \begin{align}
        (f +_q g) +_p h & =\, \diagp{X};((f +_q g) \oplus h);\codiag{Y} \tag{def. $+_p$}\\
        & =\, \diagp{X};(\diagq{X} \oplus \id{X});((f \oplus g) \oplus h);(\codiag{Y}\oplus \id{Y});\codiag{Y} \tag{def. $+_q$}\\
        & =\, \diagptilde{X};(\id{X} \oplus \diagqtilde{X});(f \oplus (g \oplus h));( \id{X}\oplus\codiag{Y} );\codiag{Y} \tag{\ref{ax:PCA1}+\ref{ax:Mon1}}\\
        & =\, \diagptilde{X};(f \oplus (g +_{\tilde{q}} h));\codiag{Y} \tag{def. $+_{\tilde{q}}$}\\
        & = f +_{\tilde{p}} (g +_{\tilde{q}} h). \tag{def. $+_{\tilde{p}}$}
    \end{align}



    We pass now to the equations in \eqref{eq:enr}. For every $e\colon Z \to X$ and $f,g\colon X \to Y$, we have:  
    \begin{align}
        e;(f +_p g) & = e;\,\diagp{X};(f \oplus g);\codiag{Y} \tag{def. $+_p$}\\
        & =\, \diagp{X};(e \oplus e);(f \oplus g);\codiag{Y} \tag{\ref{eq:nat copca1}}\\
        & =\, \diagp{X};((e;f) \oplus (e;g));\codiag{Y} \tag{SMC}\\
        & = (e;f) +_p (e;g). \tag{def. $+_{p}$}
    \end{align}
The second equation is proved similarly. The third equation is proved below for $f\colon Z\to X$.
    \begin{align}
        f;\star_{X,Y} & = f;\,\bangp{X};\cobang{Y} \tag{def. $\star$}\\
        & =\, \bangp{Z};\cobang{Y} \tag{\ref{eq:nat copca2}}\\
        & = \star_{Z,Y}. \tag{def. $\star$}
    \end{align}
\end{proof}

\begin{proof}[Proof of Lemma \ref{lemma:coproducts}]
The first two points follow immediately by the Fox theorem \cite{fox1976coalgebras}. Naturality of $\,\bangp{X}$ implies that $0$ is also a terminal object, hence a zero object. 
  To prove that axiom \eqref{eq:delta} holds, consider the following equalities:
    \begin{align}
        \iota_1;\pi_1 & = \Irunit{X_1};(\id{X_1}\oplus \cobang{X_2});(\id{X_1}\oplus \bang{X_1});\runit{X_1} \tag{def. of $\iota_1$ and $\pi_1$}\\
        & = \Irunit{X_1};(\id{X_1}\oplus (\cobang{X_2};\bang{X_1}));\runit{X_1} \tag{SMC}\\
        & = \Irunit{X_1};(\id{X_1}\oplus \id{0});\runit{X_1} \tag{\ref{eq:nat copca2} + \ref{ax:coh3}}\\
        & = \id{X_1} \tag{SMC}
    \end{align}
    Similarly, one can prove that $\iota_2;\pi_2 = \id{X_2}$. Now, for $i\neq j$, we have: 
    \begin{align}
        \iota_1;\pi_2 & = \Irunit{X_1};(\id{X_1}\oplus \cobang{X_2});(\bang{X_1}\oplus \id{X_1});\lunit{X_1} \tag{def. of $\iota_1$ and $\pi_2$}\\
        & = \Irunit{X_1};((\id{X_1};\bang{X_1})\oplus (\cobang{X_2};\id{X_1}));\lunit{X_1} \tag{SMC}\\
        & = \Irunit{X_1};(\bang{X_1}\oplus \cobang{X_2});\lunit{X_1} \tag{SMC}\\
        & = \Irunit{X_1};(\,\bangp{X_1}\oplus \id{0});(\id{0}\oplus \cobang{X_2});\lunit{X_1}\tag{SMC}\\  
        & = \bang{X_1};\Irunit{0};\lunit{0};\cobang{X_2}    \tag{Nat. $\Irunit{},\lunit{}$}     \\
        & = \bang{X_1};\cobang{X_2} \tag{$\runit{0}=\lunit{0}$} \\
        & =\star_{X_1,X_2} \tag{Lemma~\ref{lemma:enrichmentpca}}
    \end{align}
    and similarly $\iota_2;\pi_1 = \bang{X_2};\cobang{X_1}$.
\end{proof}

\begin{proof}[Proof of Lemma \ref{lemma:diagpq}] 
We prove the statement only for $p,q\in(0,1)$ such that $p+q<1$, the other cases are similar and left to the reader.
Point \eqref{lemma:diagpq1} is proved by the following derivation.

    %\begin{align}
     %   \langle f_1,f_2\rangle_{p,q};\pi_1 & =\, \diagp{X};(\id{X}\oplus \diaggen{\frac{q}{1-p}}{X}); (f_1\oplus (f_2 \oplus \bang{X}));(\id{Y_1}\oplus \runit{Y_2});(\id{Y_1}\oplus \bang{Y_1});\runit{Y_1} \notag\\
    %    &=\, \diagp{X};(\id{X}\oplus \diaggen{\frac{q}{1-p}}{X}); (f_1\oplus (f_2 \oplus \bang{X}));(\id{Y_1}\oplus \bang{Y_2\oplus 0});\runit{Y_1} \tag{\ref{eq:nat copca2}}\\
    %    &=\, \diagp{X};(\id{X}\oplus \diaggen{\frac{q}{1-p}}{X});(f_1 \oplus \bang{X\oplus X}); \runit{Y_1} \tag{\ref{eq:nat copca2}}\\
    %    &=\, \diagp{X};(\id{X}\oplus \diaggen{\frac{q}{1-p}}{X});(f_1 \oplus ((\bang{X} \oplus \bang{X});\Ilunit{0})); \runit{Y_1} \tag{\ref{eq:coherence cobang}}\\
    %    &=\, \diagp{X};(f_1 \oplus \bang{X});\runit{Y_1} \tag{\ref{eq:star I = id I} + \ref{eq:nat copca1}}\\
    %    &= \,\diagp{X};(f_1 \oplus \bang{X});\codiag{Y_1} \tag{\ref{ax:Mon2}}\\
    %    & = p\cdot f_1. \notag TEMP \diaggen{p,q\,}{X}
    %\end{align}
    \begin{align}
        \diaggen{p,q\,}{X};\pi_1 
        &=\, \diaggen{p,q\,}{X};(\id{X}\oplus \bang{X});\runit{X} \tag{\eqref{eq:defpi}}\\
         & =\, \diagp{X};(\id{X}\oplus \diaggen{\frac{q}{1-p}}{X});(\id{X}\oplus (\id{X} \oplus\bang{X}));(\id{X}\oplus \runit{X});(\id{X}\oplus \bang{X});\runit{X} \tag{\eqref{diagpq generalizzata}}\\
         &= \, \diagp{X};(\id{X}\oplus \diaggen{\frac{q}{1-p}}{X});(\id{X}\oplus (\id{X} \oplus\bang{X}));(\id{X}\oplus (\runit{X};\,\bangp{X}));\runit{X} \tag{SMC}\\
         &= \, \diagp{X};(\id{X}\oplus \diaggen{\frac{q}{1-p}}{X});(\id{X}\oplus (\id{X} \oplus\bang{X}));(\id{X}\oplus ((\,\bangp{X}\oplus \id{0});\runit{0}));\runit{X} \tag{Nat. $\runit{}$}\\
         &= \, \diagp{X};(\id{X}\oplus \diaggen{\frac{q}{1-p}}{X});(\id{X}\oplus (\id{X} \oplus\bang{X}));(\id{X}\oplus ((\,\bangp{X}\oplus\, \bangp{0});\runit{0}));\runit{X} \tag{\ref{ax:coh3}}\\
        & =\, \diagp{X};(\id{X}\oplus \diaggen{\frac{q}{1-p}}{X});(\id{X}\oplus (\id{X} \oplus\bang{X}));(\id{X}\oplus \bang{X\oplus 0});\runit{X} \tag{\ref{eq:coherence cobang}}\\
        & =\, \diagp{X};(\id{X}\oplus \diaggen{\frac{q}{1-p}}{X});(\id{X}\oplus ((\id{X} \oplus\bang{X});\bang{X\oplus 0}));\runit{X} \tag{SMC}\\
        & =\, \diagp{X};(\id{X}\oplus \diaggen{\frac{q}{1-p}}{X});(\id{X}\oplus \bang{X\oplus X}); \runit{X} \tag{\ref{eq:nat copca2}}\\   
        & =\, \diagp{X};(\id{X}\oplus \diaggen{\frac{q}{1-p}}{X});(\id{X}\oplus ((\bang{X} \oplus \bang{X});\Ilunit{0})); \runit{X} \tag{\ref{eq:coherence cobang}}\\
        & =\, \diagp{X};(\id{X}\oplus (\diaggen{\frac{q}{1-p}}{X};(\bang{X} \oplus \bang{X});\Ilunit{0})); \runit{X} \tag{SMC}\\
        & =\, \diagp{X};(\id{X}\oplus (\,\bangp{X};\diaggen{\frac{q}{1-p}}{0};\Ilunit{0})); \runit{X} \tag{\ref{eq:nat copca1}}\\
        & =\, \diagp{X};(\id{X}\oplus\, \bang{X});\runit{X} \tag{\ref{eq:star I = id I} }\\
        & = \,\diagp{X};(\id{X}\oplus \bang{X});(\id{X}\oplus\cobang{X});\codiag{X} \tag{\ref{ax:Mon2}}\\
        & =\, \diagp{X};(\id{X}\oplus \bang{X};\cobang{X});\codiag{X} \tag{SMC}\\
        & =\, \diagp{X};(\id{X}\oplus \star_{X,X});\codiag{X} \tag{\eqref{eq:enrichment}}\\
        & = \id{X}+_p \star \tag{\eqref{eq:enrichment}}\\
        & = p\cdot \id{X}. \tag{def. of $p\cdot-$}
    \end{align}
    Point \eqref{lemma:diagpq2} is proved similarly. Point~\eqref{lemma:diagpq4} is proved by the following derivation.
    \begin{align}
         \diaggen{p,q\,}{X};(f_1\oplus f_2);(\id{Y_1}\oplus \bang{Y_1});\runit{Y_1} 
        & =\, \diaggen{p,q\,}{X};(\id{X}\oplus \bang{X}); (f_1 \oplus \id{0});\runit{Y_1} \tag{\ref{eq:nat copca2}}\\
        & =\, \diaggen{p,q\,}{X};(\id{X}\oplus \bang{X});\runit{X};f_1 \tag{Nat. $\rho$}\\
        & =\, (p\cdot\id{X});f_1 \tag{Lemma~\ref{lemma:diagpq}.\ref{lemma:diagpq1}}\\
        & = p\cdot f_1 \text{.} \tag{pca-enrichment}
    \end{align}
    Similarly, one can prove point~\eqref{lemma:diagpq5}. 
    \end{proof}

\begin{proof}[Proof of Lemma \ref{lemma:diagpq3}]
By Proposition~\ref{lemma: copca objects in convbicat}, every object in $\Cat{C}$ is equipped with natural and coherent monoid and co-pca structures. By the last two items of Lemma \ref{lemma:diagpq}, $\diaggen{p,q\,}{Z};(f_1\oplus f_2); \pi_1= p\cdot f_1$ and $\diaggen{p,q\,}{Z};(f_1\oplus f_2); \pi_2= q\cdot f_2$. Since $\Cat{C}$ has convex products, the statement holds by uniqueness.
\end{proof}

